\chapter{Referencial teórico}
\label{cap:referencial}

Este capítulo apresenta os conceitos usados na avaliação: o modelo de qualidade de produto da família SQuaRE, o processo de avaliação, o método GQM, o modelo de informação do PSM e as medidas de referência adotadas na especificação.

\section{Qualidade de produto e a família SQuaRE}
\label{sec:square}

A série ISO/IEC 25000, conhecida como SQuaRE, organiza em divisões os requisitos, os modelos, as medidas e o processo de avaliação da qualidade de sistemas e software \cite{iso25000}. Três divisões são usadas diretamente neste trabalho: o modelo de qualidade (ISO/IEC 25010), as medidas de qualidade do produto e em uso (ISO/IEC 25023 e ISO/IEC 25022) e o processo de avaliação (ISO/IEC 25040).

A ISO/IEC 25010:2023 substitui a versão de 2011 e organiza a qualidade do produto em nove características, cada uma subdividida em subcaracterísticas \cite{iso25010_2023}. Como parte da literatura e do material da disciplina ainda usa a nomenclatura de 2011, a \autoref{tab:iso25010} apresenta a correspondência. Neste texto, as características renomeadas aparecem com o nome de 2023 seguido do de 2011 entre parênteses. A satisfação, que integrava o modelo de qualidade em uso na versão de 2011 \cite{iso25010_2011}, não faz parte do modelo de produto de 2023; quando usada, é tratada como medida complementar de qualidade em uso \cite{iso25022}.

\begin{table}[htb]
	\caption{Características da ISO/IEC 25010:2023 e correspondência com a versão de 2011}
	\label{tab:iso25010}
	\centering
	\footnotesize
	\begin{tabularx}{\textwidth}{L{3.6cm} L{2.7cm} Y}
		\toprule
		\textbf{ISO/IEC 25010:2023} & \textbf{ISO/IEC 25010:2011} & \textbf{Subcaracterísticas (2023)} \\
		\midrule
		Adequação Funcional & Adequação Funcional & Completude, correção e adequação (pertinência) funcional \\
		Eficiência de Desempenho & Eficiência de Desempenho & Comportamento temporal, utilização de recursos, capacidade \\
		Compatibilidade & Compatibilidade & Coexistência, interoperabilidade \\
		Capacidade de Interação & Usabilidade & Reconhecimento de adequação, apreensibilidade, operabilidade, proteção contra erros do usuário, engajamento, inclusividade, assistência ao usuário, autodescritividade \\
		Confiabilidade & Confiabilidade & Isenção de falhas, disponibilidade, tolerância a falhas, recuperabilidade \\
		Segurança & Segurança & Confidencialidade, integridade, não repúdio, responsabilização, autenticidade, resistência \\
		Manutenibilidade & Manutenibilidade & Modularidade, reusabilidade, analisabilidade, modificabilidade, testabilidade \\
		Flexibilidade & Portabilidade & Adaptabilidade, escalabilidade, instalabilidade, substituibilidade \\
		Segurança Operacional & (não existia) & Restrição operacional, identificação de riscos, falha segura, alerta de perigo, integração segura \\
		\bottomrule
	\end{tabularx}
	\fonte{Elaborado pela equipe, com base em \citeonline{iso25010_2011} e \citeonline{iso25010_2023}.}
\end{table}

\section{Processo de avaliação de produto}
\label{sec:processo}

A ISO/IEC 25040 descreve a avaliação de produto como um processo em atividades sucessivas \cite{iso25040}. No material da disciplina, ele é apresentado em quatro fases (\autoref{fig:processo}) \cite{ramos2026processo}: (1) estabelecer os requisitos da avaliação, definindo o propósito, o tipo de produto e o modelo de qualidade; (2) especificar a avaliação, selecionando métricas, níveis de pontuação e critérios de julgamento; (3) projetar a avaliação, produzindo o plano com método, instruções, recursos e cronograma; e (4) executar a avaliação, obtendo as medidas, comparando-as com os critérios e julgando os resultados. O nível de detalhe da especificação deve ser tal que a avaliação seja repetível e reprodutível.

\begin{figure}[htb]
	\centering
	\caption{Fases do processo de avaliação de produto e entregas do projeto}
	\label{fig:processo}
	\begin{tikzpicture}[
		fase/.style={draw=blue!60!black, fill=blue!8, rounded corners, align=center,
			minimum width=2.9cm, minimum height=1.25cm, font=\footnotesize},
		seta/.style={-{Stealth[length=2.5mm]}, thick, blue!60!black},
		ent/.style={font=\scriptsize\itshape, align=center}]
		\node[fase] (f1) {\textbf{Fase 1}\\Estabelecer os\\requisitos};
		\node[fase, right=0.55cm of f1] (f2) {\textbf{Fase 2}\\Especificar a\\avaliação};
		\node[fase, right=0.55cm of f2] (f3) {\textbf{Fase 3}\\Projetar a\\avaliação};
		\node[fase, right=0.55cm of f3] (f4) {\textbf{Fase 4}\\Executar a\\avaliação};
		\draw[seta] (f1) -- (f2);
		\draw[seta] (f2) -- (f3);
		\draw[seta] (f3) -- (f4);
		\draw[thick, gray] ($(f1.south west)+(0,-0.15)$) -- ($(f2.south east)+(0,-0.15)$)
			node[ent, midway, below=2pt] {Entrega 1 (este artigo)};
		\draw[thick, gray] ($(f3.south west)+(0,-0.15)$) -- ($(f4.south east)+(0,-0.15)$)
			node[ent, midway, below=2pt] {Entrega 2 (resultados e ação de melhoria)};
	\end{tikzpicture}
	\fonte{Elaborado pela equipe, com base em \citeonline{ramos2026processo} e \citeonline{ramos2026roteiro}.}
\end{figure}

\section{Método \textit{Goal-Question-Metric}}
\label{sec:gqm}

O GQM parte do princípio de que a medição deve ser orientada a objetivos \cite{basili1994}. A definição ocorre de cima para baixo em três níveis: o \textbf{nível conceitual} (objetivo), o \textbf{nível operacional} (questões que caracterizam o objeto em relação ao foco de qualidade) e o \textbf{nível quantitativo} (métricas que respondem às questões); a interpretação segue o caminho inverso, das medidas ao objetivo \cite{ramos2026gqm}. Cada objetivo é escrito em um \textit{template} com cinco elementos: o objeto analisado, o propósito, o foco de qualidade, o ponto de vista e o contexto \cite{basili1994}. Para cada questão formula-se uma hipótese, isto é, a resposta esperada antes da medição, que depois apoia a interpretação dos resultados. Boas métricas são simples, objetivas, fáceis de coletar, robustas e válidas, ou seja, medem o que pretendem medir \cite{ramos2026gqm}.

\section{Medição prática de software e sistemas}
\label{sec:psm}

O PSM é uma abordagem de medição dirigida por necessidades de informação, com forte ligação à gestão de projetos \cite{mcgarry2002}. Seu modelo de informação, também adotado pela ISO/IEC/IEEE 15939 \cite{iso15939}, liga os atributos das entidades medidas a medidas básicas (obtidas por um método de medição), a medidas derivadas (obtidas por uma função de medição) e a indicadores (obtidos por um modelo de análise associado a critérios de decisão), que respondem a uma necessidade de informação \cite{ramos2026psm}. O PSM agrupa as necessidades de informação em categorias, entre elas calendário e progresso, recursos e custos e desempenho dos processos \cite{mcgarry2002}. O \textit{framework} para desenvolvimento contínuo e iterativo (CID) adapta essas categorias a equipes que trabalham em ciclos curtos e especifica medidas como \textit{burndown}, comprometido \textit{versus} concluído e tempo de ciclo \cite{psmcid1,psmcid2}. Essas categorias e medidas são usadas na gestão deste projeto (\autoref{sec:gestao}).

\section{Medidas de referência}
\label{sec:medidas-referencia}

As métricas da Fase 2 adotam, sempre que possível, medidas e limites reconhecidos. A ISO/IEC 25023 define medidas de produto como cobertura e correção funcional na forma $X = 1 - A/B$ \cite{iso25023}, e a ISO/IEC 25022 define medidas de eficácia, eficiência e satisfação em uso \cite{iso25022}. A taxa de conclusão de tarefas tem média de referência de 78\% em estudos de usabilidade \cite{sauro2011}. O \textit{System Usability Scale} (SUS) é um questionário de dez itens que produz uma pontuação de 0 a 100 \cite{brooke1996}, cuja média de referência é 68 \cite{sauro2016}; a \textit{Single Ease Question} (SEQ) mede a facilidade percebida de uma tarefa em um único item \cite{sauro2016}. No desempenho do lado do cliente, o \textit{Largest Contentful Paint} (LCP) é considerado bom até 2,5~s e ruim acima de 4,0~s \cite{webvitals}, e 1,0~s é o limite para que o usuário mantenha o fluxo de pensamento \cite{nielsen1993}; essas medidas podem ser obtidas com a ferramenta Lighthouse \cite{lighthouse}. A acessibilidade é verificada pelos critérios de nível A e AA das WCAG~2.2 \cite{wcag22}.
